\documentclass[11pt]{article}
\usepackage[T1]{fontenc}
\usepackage{lmodern,amsmath,amssymb,amsthm,mathtools,microtype}
\usepackage[margin=1.02in]{geometry}
\usepackage{xcolor}
\definecolor{links}{HTML}{234F66}
\usepackage[colorlinks=true,allcolors=links]{hyperref}
\usepackage{enumitem}
\hypersetup{pdftitle={The joint spectrum of small-time Li heat decay},pdfsubject={Hausdorff and packing dimensions of small-time decay and approximation exponents}}
\setlist{nosep,topsep=4pt}
\newtheorem{theorem}{Theorem}[section]
\newtheorem{lemma}[theorem]{Lemma}
\newtheorem{corollary}[theorem]{Corollary}
\newcommand{\HH}{\dim_{\mathrm H}}
\newcommand{\TT}{\mathbb T}
\newcommand{\dist}{\operatorname{dist}}
\newcommand{\eps}{\varepsilon}
\newcommand{\diam}{\operatorname{diam}}
\newcommand{\dd}{\,\mathrm d}
\newcommand{\E}{\mathbb E}
\newcommand{\ro}{\rho}
\newcommand{\RR}{\mathbb R}
\newcommand{\n}{\mathbb N}
\newcommand{\D}{\mathcal D}
\newcommand{\bA}{\mathcal A}
\newcommand{\B}{\mathcal B}
\newcommand{\T}{\mathcal T}
\newcommand{\upexp}{\overline\omega}
\newcommand{\lowexp}{\underline\omega}
\newcommand{\ranked}[1]{\textup{\texttt{#1}}}
\title{The joint spectrum of small-time Li heat decay\\[3pt]
\large With an application to best Egyptian-fraction approximation}
\author{Research note}
\date{November 27, 2026}
\begin{document}
\maketitle
\begin{abstract}
At a fixed point of the circle, a Li heat density can decay at different rates along different sequences of times. We determine the Hausdorff dimension of every set with prescribed lower and upper rates. In the normalization by $\log(1/t)\log\log(1/t)$, every pair $0\le h\le k\le 1/\log2$ occurs in every open arc. For $h>0$ its dimension is $1/(e^{1/h}-1)$, independently of $k$. For $h=0$ its Hausdorff dimension is zero. Every joint level has packing dimension one. This answers the joint-decay question in the ProveIt article using its estimates for the Dirichlet Li kernel. An alternating construction makes selected scales resistant to approximation by \emph{all} short sums of jump angles. The same argument gives an unconditional joint spectrum for the errors of best strict Egyptian-fraction underapproximations with distinct denominators.
\end{abstract}

\section{The joint decay question}
Fix a primitive Dirichlet character $\chi$, allowing the conductor-one character. We use the critical-line hypothesis for its nontrivial zeros in identifying the arithmetic Li coefficients with the following jump process. Write $\ro=1/2+i\gamma$, set $e^{i\theta_\gamma}=1-1/\ro$, and choose $\theta_\gamma=2\arctan(1/(2\gamma))$ with $\theta_0=\pi$. On $\TT=\RR/(2\pi\mathbb Z)$ let
\[
 \lambda_\chi(n)=\lim_{U\to\infty}\sum_{|\gamma|\le U}(1-e^{in\theta_\gamma}),
 \qquad p_t(x)=\frac1{2\pi}\sum_{n\in\mathbb Z}e^{-t\lambda_\chi(n)-inx}.
\]
The zeros are repeated with multiplicity and the truncation in the ordinate is symmetric. The resulting semigroup has a strictly positive density. Put
\[
 L=\log(1/t),\qquad R_t(x)=\frac{-\log p_t(x)}{L\log L},\qquad
 F(x)=\liminf_{t\downarrow0}R_t(x),\quad G(x)=\limsup_{t\downarrow0}R_t(x).
\]
The single-rate theorem in \cite{repoLi} determines the dimensions of the level sets of $F$, the almost-everywhere value of both rates and a dense set realizing their extreme separation. It asks for the joint law of dimensions in Question \ranked{li:q:decay}.
\begin{theorem}\label{th:main}
For every open nonempty arc $I\subset\TT$ and every $0\le h\le k\le1/\log2$, the set
\[
 E_I(h,k)=\{x\in I:F(x)=h,\ G(x)=k\}
\]
is nonempty and
\[
 \HH E_I(h,k)=\begin{cases}
 (e^{1/h}-1)^{-1},&h>0,\\ 0,&h=0.
 \end{cases}
\]
All other pairs have empty level sets. At every point of $E_I(h,k)$ the set of subsequential limits of $R_t(x)$ as $t\downarrow0$ is $[h,k]$.
\end{theorem}
For example, the points with constant decay rate $1/\log3$ and the points oscillating between $1/\log3$ and $1/\log2$ both have dimension $1/2$ in every arc. Taking the union of $E_I(h,1/\log2)$ over $h<1/\log2$ arbitrarily close to $1/\log2$ shows that points at which $R_t$ has no limit have full Hausdorff dimension in every arc. That set has Lebesgue measure zero by the almost-everywhere theorem of \cite{repoRec}.

\section{From heat decay to the number of jumps}
Let $S_m$ consist of all sums of at most $m$ zero angles, modulo $2\pi$, repetitions allowed, and include $0$. Since the only possible accumulation point of the distinct angles is $0$, each $S_m$ is compact. Write $\delta_m(x)=\dist(x,S_m)$ and
\[
 W_m(x)=\frac1m\log\log(1/\delta_m(x)),\quad
 \lowexp(x)=\liminf_m W_m(x),\quad \upexp(x)=\limsup_m W_m(x).
\]
Only large $m$ matter. If $\delta_m(x)=0$, set $W_m(x)=+\infty$. Reciprocals of $+\infty$ are zero.

Here are the heat estimates from \cite{repoLi,repoRec} used below, with the required uniformities. Let $a(t)\sim tL/2$. One has $\|p_t\|_\infty\ll1/a$ and $p_t(z)\gg 1/a$ on $|z|\le a$. If $m=O(\log L)$ and $\delta_m(x)\le t$, then
\begin{equation}\label{near}
 p_t(x)\gg \frac{e^{-tm}t^m}{a\,m!}.
\end{equation}
Indeed, remove one unit-rate Poisson process at each distinct angle in a representation using $j\le m$ angles. The remaining density differs from $p_t$ in sup norm by $O(a^{-1}tm e^{2tm})$, by Fourier inversion and the absolute Fourier sum $O(1/a)$. The required counts of the removed processes have joint probability at least $e^{-tm}t^j/j!$. Their residual displacement from $x$ is at most $t\le a$.

Conversely, for fixed $c>0$ and $0<b<1$, uniformly in $x$ as $L\to\infty$,
\begin{equation}\label{far}
 m=\lfloor c\log L\rfloor,\quad \delta_m(x)\ge e^{-L^b}
 \quad\Longrightarrow\quad R_t(x)\ge c+o(1).
\end{equation}
For the estimate, split the jumps at size $e^{-L^{b'}}$ with $b<b'<1$. Finite quadratic jump mass gives total large-jump rate $O(e^{2L^{b'}})$. The Poisson probability of more than $m$ such jumps is at most $\exp(-(m+1)(L-O(L^{b'})))$. On the event of fewer than $m+1$, the remaining displacement has size at least $e^{-L^b}$. The small-jump density there is $O(a^{-1}\exp(-\tfrac12e^{L^{b'}-L^b}+o(1)))$, by splitting the time in half and applying the exponential moment bound to either half. This is Lemma \ranked{li:lem:decay} of \cite{repoLi}. Both estimates use multiplicities through the kernel. The sets $S_m$ depend only on the available jump locations.

\begin{lemma}\label{lem:transfer}
For every $x$, $F(x)=1/\upexp(x)$ and $G(x)=1/\lowexp(x)$.
\end{lemma}
\begin{proof}
For $0<A<\lowexp(x)$, put $m=\lceil (\log L)/A\rceil$. Eventually $\delta_m(x)\le\exp(-e^{Am})\le t$. By \eqref{near}, $R_t(x)\le m/\log L+o(1)$, so $G(x)\le1/A$. For $B>\lowexp(x)$ finite there are arbitrarily large $m$ with $\delta_m(x)\ge\exp(-e^{Bm})$. Choose $0<c<1/B$, then $cB<b<1$. Along the times with $L=\exp((m+1/2)/c)$, one has $\lfloor c\log L\rfloor=m$ and $e^{Bm}\le L^b$. Equation \eqref{far} gives $G(x)\ge c$. These two inequalities give the identity for $G$, including $\lowexp=\infty$, since the sup norm bound gives $F\ge0$.

For $0<A<\upexp(x)$ use the infinitely many indices $m$ with $\delta_m\le\exp(-e^{Am})$, taking $L=e^{Am}$ in \eqref{near}, to get $F\le1/A$. For $B>\upexp(x)$ finite, the hypothesis $\delta_m\ge\exp(-e^{Bm})$ holds eventually for every $m$. With $c<1/B$ and $m=\lfloor c\log L\rfloor$ the far estimate applies for every sufficiently large $L$, giving $F\ge c$. The zero-count approximation below ensures $\lowexp\ge\log2$.
\end{proof}

\section{Upper bounds from all competing sums}
The proof now uses two properties of the distinct jump magnitudes $b_1\ge b_2\ge\cdots>0$. First, $b_j\le C_sj^{-1/s}$ for every $s>1$. Second, for some $r_*,C>0$ every $0<r<r_*$ admits an actual positive angle in $[r-Cr^2,r]$. Indeed the one-sided Dirichlet zero count is $T\log(qT/(2\pi e))/(2\pi)+O_\chi(\log(2+T))$. A sufficiently large constant-length interval contains an ordinate for all large $T$, and $|\mathrm d\theta/\mathrm dT|=(T^2+1/4)^{-1}$. The counting upper bound gives the first property. On the circle repeated uses of a fixed small positive angle reduce any positive representative to a residual in $[0,r_*)$ in boundedly many steps. Subtracting a suitable angle then reduces each residual $u$ to at most $Cu^2$. Thus $\delta_m(x)\le e^{-c2^m}$ for all large $m$, with a positive constant. In particular $\lowexp\ge\log2$. The finite-sum centres are dense.

Let $S_m^\pm$ allow arbitrary signs and repetitions of any magnitude. Set $d_m^\pm(x)=\dist(x,S_m^\pm)$, using circle distance or ordinary distance on a bounded real interval as appropriate. Define $\upexp^\pm$ by using $d_m^\pm$ in place of $\delta_m$. We repeat the covering argument in \cite{repoLi} to make clear why the construction below controls competing approximants.
\begin{lemma}\label{lem:cover}
If $R>2$, then $\HH\{\upexp^\pm\ge\log R\}\le 1/(R-1)$. The same bound holds for ordinary distance on the real line when the magnitudes satisfy the first property above.
\end{lemma}
\begin{proof}
Fix $1/(R-1)<d<1$. Choose $p>1/d$, $s>1$ such that $1+sp<R$. Choose $0<\eta<1$ with $2\eta^d\zeta(pd)<1$. Around every ordered signed sum $y$ of length $k$ with indices $j_1,\ldots,j_k$ draw an interval of radius $v=\eta^k\prod j_i^{-p}$. The sum of the $d$th powers of their diameters is finite, bounded by $2^d\sum_k(2\eta^d\zeta(pd))^k$. Hence, except for a set of zero $d$-dimensional measure and the countably many centres, every $x$ satisfies $\dist(x,y)\ge c_x v$ for all tuples, with $c_x>0$.

For a particular tuple of length $n\le m$ order its terms by decreasing magnitude. Write $r_k$ for the distance from $x$ to the partial sum after $k$ terms, $B_k=\sum_{i\le k}\log j_i$ and $\Delta=r_n$. Before the first $k$ with $r_k\le2\Delta$, one has $r_k\le\Delta+mb_{j_{k+1}}$, so $b_{j_{k+1}}>r_k/(2m)$. It follows that
\[
 B_{k+1}\le(1+sp)B_k+s\{\log(2m C_s/c_x)+k\log(1/\eta)\}.
\]
Geometric summation gives $B_k\ll_x(1+sp)^k\log(2m)$ up to the stopping index. There $2\Delta\ge r_k\ge c_x\eta^k e^{-pB_k}$. Taking the infimum shows $d_m^\pm\ge\exp(-C_x(1+sp)^m\log(2m))$, proving $\upexp^\pm\le\log(1+sp)<\log R$ off the exceptional set. Let $d\downarrow1/(R-1)$.
\end{proof}

\section{Prescribing both approximation exponents}
For $2\le P\le Q\le\infty$, our construction yields
\begin{equation}\label{geom}
 \lowexp(x)=\log P,\qquad \upexp(x)=\log Q,
\end{equation}
with Hausdorff dimension $D_Q=1/(Q-1)$, where $D_\infty=0$. It will also give full packing dimension, written $\dim_{\mathrm P}$, even for the zero-dimensional Hausdorff levels.

\paragraph{Intervals and mass.}
Start with $J_0=[z_0,z_0+r_0]$ inside the specified open arc, with $z_0$ a sum of $s_0$ actual angles. Work in an ordinary coordinate on this interval. Write $H_m=\log(1/r_m)$ and use common lengths at each level. For each step with exponent $V_m\ge2$, put
\[
 r_{m+1}=\kappa r_m^{V_m},\qquad H_{m+1}=V_m H_m+u,\qquad u=\log(1/\kappa)>0.
\]
There are $M_m\asymp1/r_m$ actual positive angles in $[r_m/4,3r_m/4]$ that are mutually separated by at least $c_0r_m^2$. To see this, place test values in $[r_m/3,2r_m/3]$ at spacing $K r_m^2$, and use the one-sided angle property for large fixed $K$. Keep the same angles for every parent at this level. Each child of $[z,z+r_m]$ is $[z+y,z+y+r_{m+1}]$ for one of these $y$. Fix $0<\kappa<\min(c_0/4,1/8)$ and $r_0$ small enough. Children lie inside their parent and are disjoint, with gaps at least $c_0r_m^2/2$. Each point of the limit set satisfies
\begin{equation}\label{treeapprox}
 \delta_{s_0+m}(x)\le r_m,\qquad
 W_{s_0+m}(x)\ge\frac{\log H_m}{s_0+m}.
\end{equation}
The bounded initial index $s_0$ has no effect on the exponents.

Put $A_m=\prod_{i<m}M_i$, so $\log A_m=\sum_{i<m}H_i+O(m)$. We will occasionally prune an entire finite block at its end. In each retained parent at the start of such a block keep at least half the terminal descendants, and divide the parent's mass equally among those retained. Divide equally at all unpruned steps. This determines a probability measure $\mu$ on the limit set. Inside a pruned block the mass of any partial descendant is at most twice what uniform division within that block would give, since its fraction of unpruned terminals bounds its numerator and the retained fraction is at least $1/2$. Consequently every level-$m$ interval has mass at most $2^{m+1}/A_m$. This estimate includes partial levels and holds for the measure after all subsequent choices.

If $V_i\le Q<\infty$, telescoping $H_{i+1}-H_i\le(Q-1)H_i+u$ gives
\[
 \log A_m\ge D_Q H_m-O(m+H_0).
\]
Since $H_m\ge2^m H_0$, the measure satisfies $\mu(J)\ll_d |J|^d$ for every interval $J$ and every $0<d<D_Q$. Here are the bounds between construction scales. For $r_{m+1}\le\ell=|J|\le r_m$, the interval meets at most two level-$m$ parents. Their spacing between children gives at most $C(1+\ell/r_m^2)$ children in each. When $\ell\le r_m^2$, bound the mass by $C2^{m+2}/A_{m+1}\ll_d r_{m+1}^d$. When $\ell\ge r_m^2$, the same count gives $C(2^{m+2}/A_m)\ell/r_m\ll_d r_m^d\ell/r_m\le C_d\ell^d$. These are the claimed Frostman estimates. In particular they also hold for an infinite homogeneous continuation $V_i=P$ from any fixed cylinder with its own uniform probability measure and all $d<D_P$.

\paragraph{Slow stages protect against competing sums.}
Assume $P<\infty$. At the beginning of the $j$th slow stage, there are finitely many retained parents, all at a common depth. For each parent consider its \emph{auxiliary} infinite homogeneous continuation of exponent $P$, before choosing the length of the actual stage. Its measure is Frostman for every $d<D_P$. By Lemma~\ref{lem:cover}, it gives zero mass to $\{\upexp^\pm\ge\log P+1/(2j)\}$, since $1/(P e^{1/(2j)}-1)<D_P$. In particular almost every point of that auxiliary continuation satisfies, for every sufficiently large \emph{absolute} index $n$,
\[
 d_n^\pm(x)\ge 2 \tau_n,\qquad
 \tau_n=\exp\{-\exp[n(\log P+1/j)]\}.
\]
Choose a common $n_j$ larger than the current absolute depth and all preceding selected indices, so that the displayed good set at $n=n_j$ has measure at least $1/2$ under \emph{each} auxiliary law. This follows by convergence of its indicators to one and finiteness of the list of parents. Only now choose a common terminal level $m'>n_j$ along these homogeneous continuations deep enough that $r_{m'}<\tau_{n_j}$. In each parent retain all terminal cylinders meeting its good set. They make up at least half the unpruned terminal cylinders. The function $d_{n_j}^\pm$ is 1-Lipschitz, so throughout each entire retained cylinder it is at least $\tau_{n_j}$. Later choices preserve this property. Since $S_{n_j}\subseteq S_{n_j}^\pm$, all surviving points have
\begin{equation}\label{protect}
 W_{n_j}(x)\le\log P+1/j.
\end{equation}
The good-set argument examines every signed sum of length at most $n_j$, including all competitors to the branch prefixes. It does not assume the asserted spectrum for any auxiliary measure.

\paragraph{Fast stages and sparse quadratic stages.}
If $Q<\infty$ put $B_j=Q$. For $Q=\infty$, $P<\infty$, put $B_j=P+j$. If $P=Q=\infty$ use $B_j=j+2$ and omit the slow stages entirely. At stage $j$, repeat the exponent $B_j$ long enough that its terminal depth $m$ satisfies $m\ge j^4$ and
\begin{equation}\label{fast}
 \frac{\log H_m}{m+s_0}\ge\log B_j-1/j .
\end{equation}
Such a depth exists since under a fixed exponent $B_j$, the left side tends to $\log B_j$ as a function of the length of the extension. Follow this fast block by $j$ steps of exponent exactly $2$, then take the slow stage described above if $P<\infty$. Repeat forever. At every intermediate depth $m$, the number of steps belonging to the inserted quadratic blocks is $o(m)$: by the time the $j$th such block starts, the cumulative number through it is at most $j(j+1)/2$ and its starting depth is at least $j^4$.

If $P<\infty$, outside those sparse steps $V_m\ge P$. Summing $\log H_{i+1}-\log H_i\ge\log V_i$ in \eqref{treeapprox} gives $\lowexp\ge\log P$ everywhere on the limit set. Equation \eqref{protect} gives the reverse inequality. For $P=\infty$ the nonsparse $V_m$ tend to infinity, so the same sum gives $\lowexp=\infty$. Equation \eqref{fast} gives $\upexp\ge\log Q$ in all cases. For finite $Q$, the Frostman bounds imply that $\mu$ vanishes on every $\{\upexp^\pm\ge\log Q+1/l\}$, $l\in\n$, by Lemma~\ref{lem:cover}. Removing their union leaves a full-$\mu$ set where $\upexp=\log Q$. This proves the lower dimension bound $D_Q$. The upper bound follows from the same covering lemma with $R=Q$ for $Q>2$, from the ambient dimension for $Q=2$, and by sending $R\to\infty$ for $Q=\infty$. For $Q=\infty$ no points of the limit set need be removed to enforce \eqref{geom}.

\paragraph{The packing dimension.}
At the endpoint $v=m+j$ of each quadratic block, the formula $H_{m+i}=2^i(H_m+u)-u$ yields
\[
 \sum_{i=0}^{j-1}H_{m+i}=(1-2^{-j})H_v-O(ju).
\]
It follows that $2^{v+1}/A_v\le\exp[-(1-2^{-j}-o(1))H_v]$, since $v/H_v\to0$. An interval of length $r_v$ meets at most two level-$v$ cylinders. Thus any subset of positive $\mu$-outer measure requires at least $r_v^{-(1-o(1))}$ intervals of length $r_v$ in a cover as $j\to\infty$, and has upper box dimension one. In a countable cover of any full-$\mu$ subset, at least one member has positive outer measure. The characterization of packing dimension by countably stabilized upper box dimension \cite{falconer} gives packing dimension one for every joint level in every open arc. These scale bounds also show that $\mu$ is atomless even when $Q=\infty$.

\begin{corollary}\label{cor:pack}
All the nonempty joint levels of Theorem~\ref{th:main} have full packing dimension in every open arc, including $F=G=0$. Every level contains uncountably many points outside the finite sums of zero angles.
\end{corollary}
\begin{proof}[Proof of the heat-kernel conclusions]
Take $P=e^{1/k}$, $Q=e^{1/h}$, using $e^{1/0}=\infty$. Lemma~\ref{lem:transfer} identifies the desired level with \eqref{geom}. The uniform approximation bound forces $0\le F\le G\le1/\log2$. Continuity in $t>0$ follows from strict positivity and locally uniform convergence of the Fourier series on compact time intervals. A continuous real function on $(0,t_0)$ with finite lower and upper limits has every intermediate value as a subsequential limit at zero. Finally the finite sums are countable and the constructed measures are atomless.
\end{proof}

\section{Best Egyptian-fraction errors}
For an irrational $x\in(0,1)$ let
\[
 e_m(x)=\inf\left\{x-\sum_{i=1}^m\frac1{n_i}>0:
                  2\le n_1<\cdots<n_m\right\}.
\]
Thus one optimizes afresh over all denominators at each depth. The best approximants need not be nested. For irrational $x$ this infimum is positive: a convergent subsequence of $m$-term sums has limit a sum of at most $m$ unit fractions, by compactness of $\{0,1/2,1/3,\ldots\}$. No such limit equals $x$.
\begin{theorem}\label{th:egypt}
For every nonempty open interval $I\subset(0,1)$ and $2\le P\le Q\le\infty$, the irrational numbers $x\in I$ satisfying
\[
 \liminf_{m\to\infty}\frac{\log\log(1/e_m(x))}{m}=\log P,\qquad
 \limsup_{m\to\infty}\frac{\log\log(1/e_m(x))}{m}=\log Q
\]
form a set of Hausdorff dimension $1/(Q-1)$ and packing dimension one, with Hausdorff dimension zero when $Q=\infty$. These are all the possible pairs. The cluster set of the normalized errors, in the extended real line, is the entire interval $[\log P,\log Q]$. In particular the set with different lower and upper exponents has full Hausdorff dimension in every interval and Lebesgue measure zero.
\end{theorem}
\begin{proof}
Use the real-line form of Lemma~\ref{lem:cover} with $b_j=1/(j+1)$. The distance to unrestricted signed sums satisfies $d_m^\pm(x)\le e_m(x)$. The strict greedy algorithm with $n=\lfloor1/r\rfloor+1$ sends a residual $r$ to $0<r-1/n\le r^2/(1+r)<r^2$. For irrational $r$, the residual is positive. The greedy denominators increase strictly since $r<1/(n-1)$ and $r-1/n<1/(n(n-1))\le1/n$. This supplies $e_m(x)\le C_x^{2^m}$ with $0<C_x<1$, as well as dense finite-sum centres with distinct denominators. It also gives the required quadratic spacing of available fractions at small scales. Begin with a finite-sum centre in $I$ and choose $r_0$ so small that every first-step denominator exceeds those already used. The child shifts lie in $[r_m/4,3r_m/4]$ and $3r_{m+1}/4<r_m/4$, so the denominator increases strictly down each branch. Each limit point exceeds every prefix by a positive amount at most $r_m$. Consequently \eqref{treeapprox} holds for $e_m$ with the same index shift. At slow stages use the 1-Lipschitz signed distance as before, which gives $e_{n_j}\ge d_{n_j}^\pm\ge\tau_{n_j}$. This avoids any assumption about continuity of the one-sided best error. All the mass and dimension arguments now apply. Discarding the countably many rational points preserves the full constructed measure also when $Q=\infty$. Finally the union of the exceptional sets $\{\upexp^\pm\ge\log(2+1/j)\}$ has Lebesgue measure zero by Lemma~\ref{lem:cover}. Letting $Q>2$ decrease to $2$ with $P=2$ proves the full-dimension claim for oscillation. To see that the cluster interval has no gaps, write $w_n=n^{-1}\log\log(1/e_n)$. Adding a sufficiently small unused unit fraction to any strict underapproximation shows that $e_{n+1}\le e_n$. Hence $w_{n+1}\ge n w_n/(n+1)$. For each finite $c$ strictly between the lower and upper exponents there are arbitrarily late downward crossings $w_n\ge c>w_{n+1}$, at which $c\le w_n<c(1+1/n)$.
\end{proof}
For $P=2$, $Q=3$, the residual has doubly exponential rates ranging between the exponents $\log2$ and $\log3$ in the sense of the two displayed limits; the targets have Hausdorff dimension $1/2$ and packing dimension one. This statement concerns the \emph{best} approximation error among all $m$-term sums. Kova\v c studies when successive maximizers are eventually nested \cite{kovac}; Kova\v c and Tang establish eventual greediness at rational targets and construct an irrational example with unique greedy best underapproximations \cite{kt}. Wu's Hausdorff dimension theorem concerns the growth of successive digits of the canonical Sylvester expansion \cite{wu}. A priori none of those statements controls errors of competing $m$-term sums at the points selected here.

\section*{Scope and further questions}
The signed-sum covering lemma and the construction of sets with prescribed approximation exponents are real-variable arguments. Theorem~\ref{th:egypt} uses no hypothesis on zeros. For the arithmetic heat result, the input is precisely the circle kernel and the near- and far-sum estimates in Section~2. The proof uses the critical-line hypothesis in identifying the Dirichlet Li exponent and obtaining its analytic estimates, following \cite{repoLi,repoRec,lagarias}. It requires neither simplicity of zeros nor arithmetic independence of their ordinates. The argument resolves the question stated in \cite{repoLi}; it provides no independent verification of the repository's critical-line theorem. Picard's small-time theory \cite{picard} relates jump reachability to the size of a heat density. The statement proved here specifies both normalized decay exponents simultaneously and quantifies the size of each level.

The formula for Hausdorff dimension determines a critical exponent, but does not give the Hausdorff measure in the critical dimension when $0<h<1/\log2$. Determining suitable gauge functions at that endpoint would require finer control of the distribution of short sums. A second question is whether imposing the same pair of rates for two different Dirichlet characters changes the joint dimension; the construction above uses actual jumps from one specified process at every step.

\begin{thebibliography}{10}
\bibitem{repoLi} V. Reshetnikov, \emph{ProveIt, Generalized Riemann Hypothesis, Article, Li coefficients}, revision \texttt{fc7d763d}. See the \href{https://github.com/VladimirReshetnikov/ProveIt/blob/fc7d763da930f3231ecb3bafe0f49f4e0d4bb9a4/Analysis/GeneralizedRiemannHypothesis/Article/sections/li.tex\#L2301-L2391}{separation, decay and marginal-spectrum results} and the \href{https://github.com/VladimirReshetnikov/ProveIt/blob/fc7d763da930f3231ecb3bafe0f49f4e0d4bb9a4/Analysis/GeneralizedRiemannHypothesis/Article/sections/li.tex\#L2786-L2791}{joint-decay question}.
\bibitem{repoRec} V. Reshetnikov, \emph{ProveIt, Generalized Riemann Hypothesis, Article, Reconstruction}, same revision. See the \href{https://github.com/VladimirReshetnikov/ProveIt/blob/fc7d763da930f3231ecb3bafe0f49f4e0d4bb9a4/Analysis/GeneralizedRiemannHypothesis/Article/sections/rec.tex\#L918-L957}{oriented kernel} and the \href{https://github.com/VladimirReshetnikov/ProveIt/blob/fc7d763da930f3231ecb3bafe0f49f4e0d4bb9a4/Analysis/GeneralizedRiemannHypothesis/Article/sections/rec.tex\#L1179-L1203}{typical height and uniform lower estimate}.
\bibitem{lagarias} J. C. Lagarias, \emph{Li coefficients for automorphic $L$-functions}, Ann. Inst. Fourier \textbf{57} (2007), 1689--1740. \href{https://doi.org/10.5802/aif.2311}{DOI 10.5802/aif.2311}.
\bibitem{picard} J. Picard, \emph{Density in small time for L\'evy processes}, ESAIM Probab. Stat. \textbf{1} (1997), 357--389. \href{https://www.numdam.org/item/PS_1997__1__357_0/}{Article}.
\bibitem{falconer} K. Falconer, \emph{Fractal Geometry: Mathematical Foundations and Applications}, third ed., Wiley, 2014, Chapter 3.
\bibitem{kovac} V. Kova\v c, \emph{On eventually greedy best underapproximations by Egyptian fractions}. \href{https://arxiv.org/abs/2406.07218}{arXiv:2406.07218}.
\bibitem{kt} V. Kova\v c and Q. Tang, \emph{Eventually greedy best Egyptian underapproximations of rational numbers via optimal control}, preprint (2026). \href{https://arxiv.org/abs/2607.28387}{arXiv:2607.28387}.
\bibitem{wu} J. Wu, \emph{On the distribution of denominators in Sylvester expansions (II)}, Math. Proc. Cambridge Philos. Soc. \textbf{135} (2003), 421--430. \href{https://www.cambridge.org/core/journals/mathematical-proceedings-of-the-cambridge-philosophical-society/article/abs/on-the-distribution-of-denominators-in-sylvester-expansions-ii/65D53EA053053A44168312C42B21813E}{Article}.
\end{thebibliography}
\end{document}
